\textbf{Phần 1. Trắc nghiệm lựa chọn}
\setcounterpageref{ex}{0}

\begin{ex} %[5H6N32]
Hình nào dưới đây có hai cặp cạnh đối diện song song và bốn cạnh bằng nhau?
\choice
{Hình chữ nhật}
{Hình bình hành}
{\True Hình thoi}
{Hình thang}
\loigiai{
Hình thoi là hình có hai cặp cạnh đối diện song song và bốn cạnh bằng nhau.\\
\textbf{Đáp án của câu hỏi là C}
}
\end{ex}

\begin{ex} %[5H6N32]
Trong một hình tam giác, đoạn thẳng vuông góc kẻ từ một đỉnh xuống cạnh đối diện được gọi là
\choice
{\True Đường cao}
{Cạnh đáy}
{Bán kính}
{Đường kính}
\loigiai{
Đoạn thẳng vuông góc kẻ từ đỉnh xuống cạnh đối diện là đường cao của tam giác.\\
\textbf{Đáp án của câu hỏi là A}
}
\end{ex}

\begin{ex} %[5H6H32]
Một tờ giấy hình chữ nhật có kích thước $12\text{ cm} \times 5\text{ cm}$. Người ta cắt bỏ hai hình tam giác vuông ở hai góc có các cạnh góc vuông lần lượt là $2\text{ cm} \times 5\text{ cm}$ và $5\text{ cm} \times 5\text{ cm}$. Diện tích phần tờ giấy còn lại là
\choice
{$60\text{ cm}^2$}
{\True $42{,}5\text{ cm}^2$}
{$45\text{ cm}^2$}
{$37{,}5\text{ cm}^2$}
\loigiai{
Diện tích tờ giấy ban đầu là: $12 \times 5 = 60\text{ (cm}^2\text{)}$.\\
Diện tích hai hình tam giác vuông cắt đi là: $\dfrac{2 \times 5}{2} + \dfrac{5 \times 5}{2} = 5 + 12{,}5 = 17{,}5\text{ (cm}^2\text{)}$.\\
Diện tích phần giấy còn lại là: $60 - 17{,}5 = 42{,}5\text{ (cm}^2\text{)}$.\\
\textbf{Đáp án của câu hỏi là B}
}
\end{ex}

\begin{ex} %[5H6H32]
Một mặt hồ có dạng nửa hình tròn với bán kính đo được là $60\text{ m}$. Diện tích mặt hồ đó là
\choice
{$11\,304\text{ m}^2$}
{$188{,}4\text{ m}^2$}
{\True $5\,652\text{ m}^2$}
{$376{,}8\text{ m}^2$}
\loigiai{
Diện tích cả hình tròn bán kính $60\text{ m}$ là: $3{,}14 \times 60 \times 60 = 11\,304\text{ (m}^2\text{)}$.\\
Diện tích mặt hồ (nửa hình tròn) là: $11\,304 : 2 = 5\,652\text{ (m}^2\text{)}$.\\
\textbf{Đáp án của câu hỏi là C}
}
\end{ex}

\begin{ex} %[5H6H32]
Chu vi của miệng một chiếc bình hình tròn có bán kính $5\text{ cm}$ là
\choice
{$15{,}7\text{ cm}$}
{$78{,}5\text{ cm}$}
{$62{,}8\text{ cm}$}
{\True $31{,}4\text{ cm}$}
\loigiai{
Chu vi miệng bình là: $3{,}14 \times 2 \times 5 = 31{,}4\text{ (cm)}$.\\
\textbf{Đáp án của câu hỏi là D}
}
\end{ex}

\begin{ex} %[5H6H32]
Một hình tam giác có độ dài đáy là $8\text{ cm}$ và chiều cao tương ứng là $5\text{ cm}$. Diện tích của hình tam giác đó là
\choice
{\True $20\text{ cm}^2$}
{$40\text{ cm}^2$}
{$13\text{ cm}^2$}
{$10\text{ cm}^2$}
\loigiai{
Diện tích hình tam giác là: $\dfrac{8 \times 5}{2} = 20\text{ (cm}^2\text{)}$.\\
\textbf{Đáp án của câu hỏi là A}
}
\end{ex}

\begin{ex} %[5H6V32]
Rô-bốt có một sợi dây chun dài $15{,}85\text{ cm}$ và có thể bị đứt nếu kéo dài quá $2$ lần chiều dài ban đầu. Rô-bốt dùng sợi chun này để buộc miệng bình hình tròn có bán kính $5\text{ cm}$. Khẳng định nào sau đây đúng?
\choice
{Sợi dây chun sẽ bị đứt vì chu vi miệng bình dài hơn $2$ lần sợi dây}
{\True Sợi dây chun không bị đứt và buộc kín được miệng bình}
{Sợi dây chun quá ngắn không thể quấn quanh miệng bình}
{Sợi dây chun không cần dãn ra cũng vừa vặn miệng bình}
\loigiai{
Chu vi miệng bình là: $3{,}14 \times 2 \times 5 = 31{,}4\text{ cm}$.\\
Khi kéo dãn quanh miệng bình, chiều dài dây gấp: $31{,}4 : 15{,}85 \approx 1{,}98$ lần $< 2$ lần.\\
Do đó dây chun không bị đứt và buộc được.\\
\textbf{Đáp án của câu hỏi là B}
}
\end{ex}

\begin{ex} %[5H6V32]
Một hình bình hành có độ dài đáy $15\text{ cm}$ và chiều cao tương ứng $8\text{ cm}$. Một hình thoi có diện tích bằng diện tích hình bình hành đó và có một đường chéo dài $20\text{ cm}$. Độ dài đường chéo còn lại của hình thoi là
\choice
{$6\text{ cm}$}
{$24\text{ cm}$}
{$10\text{ cm}$}
{\True $12\text{ cm}$}
\loigiai{
Diện tích hình bình hành là: $15 \times 8 = 120\text{ (cm}^2\text{)}$.\\
Diện tích hình thoi bằng $120\text{ cm}^2$.\\
Độ dài đường chéo còn lại là: $120 \times 2 : 20 = 12\text{ (cm)}$.\\
\textbf{Đáp án của câu hỏi là D}
}
\end{ex}

\begin{ex} %[5H6V32]
Cắt ghép $4$ hình tam giác vuông có hai cạnh góc vuông là $3\text{ cm}$ và $4\text{ cm}$ để tạo thành một hình thoi. Diện tích của hình thoi đó là
\choice
{\True $24\text{ cm}^2$}
{$12\text{ cm}^2$}
{$48\text{ cm}^2$}
{$20\text{ cm}^2$}
\loigiai{
Diện tích mỗi hình tam giác vuông là: $\dfrac{3 \times 4}{2} = 6\text{ (cm}^2\text{)}$.\\
Diện tích hình thoi tạo thành là: $6 \times 4 = 24\text{ (cm}^2\text{)}$.\\
\textbf{Đáp án của câu hỏi là A}
}
\end{ex}

\begin{ex} %[5H6V32]
Một tấm bìa hình thang có đáy lớn $18\text{ cm}$, đáy bé $12\text{ cm}$. Nếu kéo dài đáy lớn thêm $4\text{ cm}$ thì diện tích tăng thêm $24\text{ cm}^2$. Diện tích tấm bìa hình thang ban đầu là
\choice
{$90\text{ cm}^2$}
{\True $180\text{ cm}^2$}
{$360\text{ cm}^2$}
{$120\text{ cm}^2$}
\loigiai{
Phần diện tích tăng thêm là tam giác có đáy $4\text{ cm}$ và cùng chiều cao với hình thang.\\
Chiều cao hình thang là: $24 \times 2 : 4 = 12\text{ (cm)}$.\\
Diện tích hình thang ban đầu: $\dfrac{(18 + 12) \times 12}{2} = 180\text{ (cm}^2\text{)}$.\\
\textbf{Đáp án của câu hỏi là B}
}
\end{ex}

\begin{ex} %[5H6V32]
Một mảnh vườn hình chữ nhật có chu vi $60\text{ m}$, chiều dài hơn chiều rộng $6\text{ m}$. Người ta đào một cái ao hình tròn bán kính $3\text{ m}$ ở giữa vườn. Diện tích phần đất còn lại là
\choice
{$216\text{ m}^2$}
{$28{,}26\text{ m}^2$}
{$197{,}74\text{ m}^2$}
{\True $187{,}74\text{ m}^2$}
\loigiai{
Nửa chu vi mảnh vườn: $60 : 2 = 30\text{ (m)}$.\\
Chiều dài mảnh vườn: $(30 + 6) : 2 = 18\text{ (m)}$.\\
Chiều rộng mảnh vườn: $30 - 18 = 12\text{ (m)}$.\\
Diện tích mảnh vườn: $18 \times 12 = 216\text{ (m}^2\text{)}$.\\
Diện tích cái ao: $3{,}14 \times 3 \times 3 = 28{,}26\text{ (m}^2\text{)}$.\\
Diện tích đất còn lại: $216 - 28{,}26 = 187{,}74\text{ (m}^2\text{)}$.\\
\textbf{Đáp án của câu hỏi là D}
}
\end{ex}

\begin{ex} %[5H6V32]
Cho hình vuông có diện tích $36\text{ cm}^2$. Hình tròn nội tiếp bên trong hình vuông tiếp xúc với cả $4$ cạnh của hình vuông có diện tích là
\choice
{$113{,}04\text{ cm}^2$}
{$56{,}52\text{ cm}^2$}
{\True $28{,}26\text{ cm}^2$}
{$9\text{ cm}^2$}
\loigiai{
Vì $6 \times 6 = 36$ nên cạnh hình vuông là $6\text{ cm}$.\\
Đường kính hình tròn bằng cạnh hình vuông là $6\text{ cm} \Rightarrow$ Bán kính $r = 3\text{ cm}$.\\
Diện tích hình tròn là: $3{,}14 \times 3 \times 3 = 28{,}26\text{ (cm}^2\text{)}$.\\
\textbf{Đáp án của câu hỏi là C}
}
\end{ex}

\textbf{Phần 2. Trắc nghiệm đúng sai}
\setcounterpageref{ex}{0}

\begin{ex} %[5H6H32]
Xét tính đúng sai của các khẳng định sau về các hình phẳng:
\choiceTF
{\True Hình thoi có hai đường chéo vuông góc với nhau và cắt nhau tại trung điểm của mỗi đường.}
{Hình thang có hai cặp cạnh đối diện song song.}
{\True Diện tích tam giác vuông bằng nửa tích độ dài hai cạnh góc vuông.}
{\True Chu vi hình tròn có bán kính $r$ tỉ lệ thuận với bán kính $r$.}
\loigiai{
\begin{itemize}
    \item Hình thoi có 2 đường chéo vuông góc tại trung điểm.
    \item Hình thang chỉ có đúng một cặp cạnh đối diện song song (hai đáy).
    \item Diện tích tam giác vuông $= \frac{a \times b}{2}$.
    \item Chu vi $C = 2 \times 3{,}14 \times r$ tỉ lệ thuận với $r$.
\end{itemize}
\begin{itemize}
    \item a) \textbf{Đúng}.
    \item b) \textbf{Sai}.
    \item c) \textbf{Đúng}.
    \item d) \textbf{Đúng}.
\end{itemize}
}
\end{ex}

\begin{ex} %[5H6V32]
Một mặt hồ có dạng nửa hình tròn bán kính $60\text{ m}$. Xét tính đúng sai của các khẳng định sau:
\choiceTF
{\True Đường kính của mặt hồ đó là $120\text{ m}$.}
{Diện tích toàn bộ hình tròn có cùng bán kính là $5\,652\text{ m}^2$.}
{\True Diện tích của mặt hồ (nửa hình tròn) là $5\,652\text{ m}^2$.}
{\True Độ dài đường bờ cong bao quanh nửa mặt hồ là $188{,}4\text{ m}$.}
\loigiai{
\begin{itemize}
    \item Đường kính: $60 \times 2 = 120\text{ m}$.
    \item Diện tích cả hình tròn: $3{,}14 \times 60 \times 60 = 11\,304\text{ m}^2 \neq 5\,652\text{ m}^2$.
    \item Diện tích nửa hình tròn: $11\,304 : 2 = 5\,652\text{ m}^2$.
    \item Độ dài bờ cong: $3{,}14 \times 60 = 188{,}4\text{ m}$.
\end{itemize}
\begin{itemize}
    \item a) \textbf{Đúng}.
    \item b) \textbf{Sai}.
    \item c) \textbf{Đúng}.
    \item d) \textbf{Đúng}.
\end{itemize}
}
\end{ex}

\textbf{Phần 3. Trả lời ngắn}
\setcounterpageref{ex}{0}

\begin{ex} %[5H6H32]
Tính diện tích của hình tam giác có độ dài đáy $8\text{ cm}$ và chiều cao $5\text{ cm}$ theo xăng-ti-mét vuông.
\shortans{20}
\loigiai{
\textbf{Đáp số:} $20$\\
$\dfrac{8 \times 5}{2} = 20\text{ (cm}^2\text{)}$.
}
\end{ex}

\begin{ex} %[5H6H32]
Mặt hồ dạng nửa hình tròn bán kính $60\text{ m}$. Tính diện tích mặt hồ theo mét vuông.
\shortans{5652}
\loigiai{
\textbf{Đáp số:} $5652$\\
$3{,}14 \times 60 \times 60 : 2 = 5\,652\text{ (m}^2\text{)}$.
}
\end{ex}

\begin{ex} %[5H6H32]
Tờ giấy hình chữ nhật $12\text{ cm} \times 5\text{ cm}$ bị cắt bỏ hai tam giác vuông ở góc có diện tích lần lượt là $5\text{ cm}^2$ và $12{,}5\text{ cm}^2$. Tính diện tích phần giấy còn lại theo xăng-ti-mét vuông. (Nhập kết quả dưới dạng số thập phân).
\shortans{42,5}
\loigiai{
\textbf{Đáp số:} $42{,}5$\\
$12 \times 5 - 5 - 12{,}5 = 42{,}5\text{ (cm}^2\text{)}$.
}
\end{ex}

\begin{ex} %[5H6V32]
Hình bình hành có diện tích $120\text{ cm}^2$. Một hình thoi có diện tích bằng diện tích hình bình hành và có một đường chéo $20\text{ cm}$. Tính độ dài đường chéo còn lại của hình thoi theo xăng-ti-mét.
\shortans{12}
\loigiai{
\textbf{Đáp số:} $12$\\
$120 \times 2 : 20 = 12\text{ (cm)}$.
}
\end{ex}

\begin{ex} %[5H6V32]
Hình thang có hai đáy $12\text{ cm}$ và $18\text{ cm}$, chiều cao $12\text{ cm}$. Tính diện tích hình thang đó theo xăng-ti-mét vuông.
\shortans{180}
\loigiai{
\textbf{Đáp số:} $180$\\
$\dfrac{(12 + 18) \times 12}{2} = 180\text{ (cm}^2\text{)}$.
}
\end{ex}

\begin{ex} %[5H6V32]
Hình vuông có diện tích $36\text{ cm}^2$. Tính diện tích của hình tròn tiếp xúc với $4$ cạnh của hình vuông theo xăng-ti-mét vuông. (Nhập kết quả dưới dạng số thập phân).
\shortans{28,26}
\loigiai{
\textbf{Đáp số:} $28{,}26$\\
Cạnh hình vuông là $6\text{ cm} \Rightarrow r = 3\text{ cm} \Rightarrow S = 3{,}14 \times 3 \times 3 = 28{,}26\text{ (cm}^2\text{)}$.
}
\end{ex}
